% !TEX program = xelatex
% !TEX encoding = UTF-8
%============================================================================
% TrajVis Studio -- Requirements Analysis Specification (Task 3)
% Course : Software Engineering -- 2026 Fall
% Group  : 1  /  Team: 泓泓说的都队
% Author : Chen Hongzhou (Student No. 1240007950), team leader
%
% Route  : Option 2, object-oriented requirements analysis (task description
%          Part A allows either Option 1, structured, or Option 2, OO).
% Layout : follows "Surface for Requirements Analysis.docx" for the cover,
%          the running header/footer and the front matter order.
% Every item still to be confirmed is marked \tbd{...} and prints in red.
% Diagrams live in figures/<name>.png; run tools/build-diagrams.sh to make
% them from the .puml sources.  A missing diagram prints a red notice
% instead of failing the build.
%============================================================================
\documentclass[12pt,a4paper]{report}

\usepackage{fontspec}
\setmainfont{Times New Roman}
\setsansfont{Helvetica}
\usepackage{xeCJK}
\setCJKmainfont{Songti SC}

\usepackage[a4paper,margin=2.5cm]{geometry}
\usepackage{graphicx}
\usepackage{booktabs}
\usepackage{array}
\usepackage{amsmath,amssymb}
\usepackage{xcolor}
\usepackage{float}
\usepackage{caption}
\usepackage{tikz}
\usetikzlibrary{positioning,arrows.meta,calc,fit,backgrounds}
\usepackage[hidelinks]{hyperref}

% Let long URLs break at slashes instead of running into the margin, and set
% them in the body font rather than the monospace one.
\urlstyle{same}
\makeatletter
\g@addto@macro\UrlBreaks{\do\/\do\-\do\.\do\_\do\:}
\makeatother

% Long identifiers and URLs should break instead of running into the margin.
\emergencystretch=3em

% Diagrams are included at exactly the width of the text block.  They are
% made squarish at the source (see the header comment in each .puml) so that
% scaling them to this width still leaves the labels at 8 pt or more; letting
% a diagram overhang into the margin instead would keep every figure out of
% the "0 overfull boxes" check that the rest of the report is held to, and a
% figure that sticks out of the text block is the first thing a marker sees.
\newcommand{\diagwidth}{\textwidth}

% --- Running header and footer -----------------------------------------
% Every page carries the running title of its chapter at the left of the
% header over a dark rule (the Surface template colours that rule 622423),
% and the page number centred in the footer.
\usepackage{fancyhdr}
% Times New Roman at 12pt needs a taller header box than the 12pt default;
% take the extra height off \topmargin so the text block does not move down.
\setlength{\headheight}{25pt}
\addtolength{\topmargin}{-10pt}
\definecolor{mustrule}{HTML}{622423}
\newcommand{\surfacehead}{%
  \fancyhead[L]{\small\nouppercase{\leftmark}}%
  \fancyfoot[C]{\small\thepage}%
  \renewcommand{\headrulewidth}{1.1pt}%
  \renewcommand{\headrule}{%
    \hbox to\headwidth{\color{mustrule}%
      \leaders\hrule height \headrulewidth\hfill}}%
}
\pagestyle{fancy}
\fancyhf{}
\surfacehead
% Chapter-opening pages use the "plain" style; give it the same furniture.
\fancypagestyle{plain}{\fancyhf{}\surfacehead}
% Header text is "Chapter N <title>", unnumbered front matter just its title.
\def\chaptermark#1{\markboth{Chapter \thechapter\enspace #1}{}}
\def\sectionmark#1{\markright{\thesection\enspace #1}}

% Chapter-based numbering, e.g. "Figure 2-1", "Table 3-2"
\renewcommand{\thefigure}{\thechapter-\arabic{figure}}
\renewcommand{\thetable}{\thechapter-\arabic{table}}

% The template names these two sections "Table of Contents" and "References"
% rather than the class defaults "Contents" and "Bibliography".
\renewcommand{\contentsname}{Table of Contents}
\renewcommand{\bibname}{References}

% --- Work markers -------------------------------------------------------
% Anything still to be settled prints in red so it cannot ship unnoticed.
% Run `grep -n '\\tbd' report.tex` before submitting; the list must be empty.
\newcommand{\tbd}[1]{{\color{red}\textbf{[TO BE CONFIRMED: #1]}}}
% Compact form for table cells, where the full wording would widen a column.
\newcommand{\tbc}[1]{{\color{red}\textbf{[TBC: #1]}}}
\newcommand{\src}[1]{{\small\itshape (Source: #1)}}

% --- Diagram slot -------------------------------------------------------
% Includes figures/<file>.png when the diagram has been built, and otherwise
% prints a red notice in its place.  This keeps the document compiling while
% the diagrams are still being drawn, and makes an unbuilt diagram visible
% rather than silently absent.
\newcommand{\diagram}[2]{%
  \IfFileExists{figures/#1.png}%
    {\includegraphics[width=\diagwidth]{figures/#1.png}}%
    {\fboxsep=10pt\fbox{\parbox[c][4.6cm][c]{0.86\textwidth}{\centering
       {\color{red}\bfseries [DIAGRAM NOT BUILT YET]}\\[8pt]
       {\ttfamily figures/#1.png}\\[8pt]
       {\small #2}}}}%
}

% --- Use-case description block -----------------------------------------
% The seven fields every use case in Chapter 2 is described by.  #1..\#7 are
% Actor, Goal, Preconditions, Main success scenario, Alternative flows,
% Postconditions and Requirements traced.
\newcommand{\ucdesc}[7]{%
  \begin{table}[H]
  \centering
  {\small
  \setlength{\tabcolsep}{4.5pt}
  \begin{tabular}{@{}>{\bfseries}p{2.9cm} >{\raggedright\arraybackslash}p{12.2cm}@{}}
    \toprule[1.5pt]
    Actor & #1 \\ \midrule
    Goal & #2 \\ \midrule
    Preconditions & #3 \\ \midrule
    Main success scenario & #4 \\ \midrule
    Alternative flows & #5 \\ \midrule
    Postconditions & #6 \\ \midrule
    Traces to & #7 \\
    \bottomrule[1.5pt]
  \end{tabular}}
  \end{table}
}

\begin{document}

%============================================================================
% COVER PAGE
%============================================================================
\begin{titlepage}
  \begin{tikzpicture}[remember picture, overlay]
    \foreach \off in {24, 26.5, 29} {
      \draw[line width=0.6pt]
        ([shift={(\off pt, -\off pt)}]current page.north west)
        rectangle
        ([shift={(-\off pt, \off pt)}]current page.south east);
    }
  \end{tikzpicture}

  \linespread{1.5}\selectfont
  \fontsize{15pt}{18pt}\selectfont
  \centering

  \vspace*{1.2cm}
  \includegraphics[height=3cm]{logo.jpeg}

  \vspace{2\baselineskip}
  {\bfseries MACAU UNIVERSITY OF SCIENCE AND\\ TECHNOLOGY\par}

  \vspace{\baselineskip}
  {\bfseries School of Computer Science and Engineering\par}
  {\bfseries Faculty of Innovation Engineering\par}

  \vspace{2\baselineskip}
  % The template writes literal "<<" and ">>".  Two adjacent "<" are a T1
  % ligature that renders as the guillemet <<, so put an empty group between
  % them to stop the ligature from forming.
  {\bfseries <{}<{}Software Project for Course Software Engineering>{}>\par}

  \vspace{2\baselineskip}
  \parbox{\textwidth}{\hspace*{1.5cm}%
    \begin{tabular}{@{}l@{: }>{\raggedright\arraybackslash}p{10.5cm}@{}}
      Homework ID  & Task3-Requirements Analysis \\
      Report Title & TrajVis Studio: Requirements Analysis
                     Specification \\
      Group Number & 1 \\
      Team Members & \begin{tabular}[t]{@{}l@{\hspace{1.2em}}l@{}}
                       Chen Hongzhou & 1240007950 \quad (team leader) \\
                       Zhou Guoxi    & 1240007947 \\
                       Xu Jing       & 1240007919
                     \end{tabular} \\
      Date         & \tbd{submission date} \\
    \end{tabular}}
\end{titlepage}

%============================================================================
% FRONT MATTER
% The template numbers the front matter in upper Roman (the cover counts as
% I, so the Abstract is II) and starts Arabic page 1 at Chapter 1.
%============================================================================
\pagenumbering{Roman}
\setcounter{page}{2}

\chapter*{Abstract}
\markboth{Abstract}{}
% \chapter* does not enter the contents, but the template's contents lists
% the Abstract, so add it by hand.  It must come after the \chapter*: that
% command starts with \clearpage, and a mark set before it would land on
% the page being closed.
\addcontentsline{toc}{chapter}{Abstract}

This document is the requirements analysis for \textbf{TrajVis Studio}, the
system proposed in Task~1 and specified in Task~2. It models the system
using object-oriented analysis, which is one of the two routes Part~A of the
task description permits.

The analysis is presented in the three stages the task description
prescribes for the object-oriented route. \textbf{Functional modelling}
(Chapter~\ref{ch:use-cases}) fixes the boundary of the system by naming its
actors and its use cases. \textbf{Static modelling}
(Chapter~\ref{ch:static}) identifies the classes the system is made of,
their attributes, their operations and the relationships between them.
\textbf{Dynamic modelling} (Chapter~\ref{ch:seq}) traces, for every use
case, the sequence of messages that carries it out.
Chapter~\ref{ch:state} adds a state diagram for the one object whose
behaviour is genuinely state-dependent, which Part~A welcomes as an extra.

Every model element is traced back to the requirement it satisfies, so that
the analysis can be checked against Task~2 one item at a time. Content that
depends on the customer assigned under Task~2 Part~A source~1, or on the
elicitation that is still outstanding, is printed in red
{\color{red}\textbf{like this}} and is not invented.

% The class already sets the running mark for these three lists, so no
% \markboth here: doing it beforehand would leak the next title onto the
% page being closed.
\tableofcontents
\addcontentsline{toc}{chapter}{Table of Contents}
\listoffigures
\addcontentsline{toc}{chapter}{List of Figures}
\listoftables
\addcontentsline{toc}{chapter}{List of Tables}

%============================================================================
\clearpage
\pagenumbering{arabic}

\chapter{Introduction}
\label{ch:intro}

\section{Purpose}
\label{sec:purpose}

This document is the Requirements Analysis Specification (RAS) for
\textbf{TrajVis Studio}, prepared by Group~1 for Task~3. Task~2 said
\emph{what} the system must do; this document says \emph{how the system is
put together} to do it.

It serves three readers. The \textbf{build team} takes the class diagram and
the sequence diagrams as the design it implements against, so that Task~4
and the code that follows are a continuation of this document rather than a
fresh start. The \textbf{assigned customer} sees the use cases as the agreed
boundary of the system and can object to a use case that is missing or
wrongly scoped before any code exists. The \textbf{course assessors} check
that every model element can be traced to a stated requirement.

\section{Notation}
\label{sec:notation}

All diagrams in this document use the Unified Modelling Language (UML), and
all of them were drawn with a diagram tool rather than by hand, as Part~B(2)
of the task description requires. In addition:

\begin{itemize}
  \item Every use case carries an identifier \texttt{UC-$n$} and every class
    an identifier of the form \texttt{C-$n$}. Requirements keep the
    identifiers they were given in Task~2, \texttt{REQ-1.1} to
    \texttt{REQ-1.6}.
  \item In the class diagram, an association end marked with a filled diamond
    is a composition and one marked with a hollow diamond is an aggregation;
    an unmarked end is a plain association.
  \item In the sequence diagrams, a solid arrow is a synchronous call, a
    dashed arrow is a return, and a bar on the lifeline is the period during
    which the participant is active.
  \item \tbd{one line explaining any other notational convention used later
    in the document, if any is introduced}
\end{itemize}

\section{Scope}
\label{sec:scope}

\subsection{In scope}
The analysis covers the whole round trip the system performs: understanding
a question written in ordinary language, retrieving and filtering the
records the question refers to, choosing and drawing the figure, describing
the result in words, recording the exchange so it can be re-run, and
measuring the system's own accuracy against a recorded benchmark.

\subsection{Out of scope}
The analysis does not cover the collection or the cleaning of the movement
data, since the system reads records that already exist. It does not cover
planning decisions, a mobile client or a public web service, all of which
Task~2 placed out of scope and which therefore have no use case here.

\section{Relation to the Earlier Deliverables}
\label{sec:relation}

Task~1 proposed the project and Task~2 stated the requirements; this
document is the third and last step before design and coding. The mapping
is direct: the six functional requirements \texttt{REQ-1.1} to
\texttt{REQ-1.6}, which came from features F1 to F6 of Task~1, become the
six use cases \texttt{UC-1} to \texttt{UC-6} of Chapter~\ref{ch:use-cases}.
Each of those use cases is then realised by the classes of
Chapter~\ref{ch:static} and by one sequence diagram in
Chapter~\ref{ch:seq}.

Task~2 left 78 items to be confirmed, most of them waiting on the customer
assigned under Part~A source~1 and on the elicitation that has not yet been
carried out. This analysis does not wait for them: the parts of the model
that are independent of the customer are complete, and the parts that are
not are marked in red and left open.

\section{Contributions}
\label{sec:contributions}

% Required by the second point of the course notice of 2026-10-08: every
% report must name what each member contributed to that report.
\tbd{name, for each of the three members, what that member contributed to
this report. The division has to be settled first; do not write it by
guessing. Once the diagrams are drawn and the use cases written, the split
is a real one: two members can take three use cases each, with the class
diagram, the sequence diagrams and the write-up divided along the same
lines as Table~\ref{tab:pairing}.}

\section{Overview}
\label{sec:overview}

Chapter~\ref{ch:use-cases} names the actors and the use cases and describes
each use case in full. Chapter~\ref{ch:static} identifies the classes and
presents the class diagram and an object diagram.
Chapter~\ref{ch:seq} gives a sequence diagram for every use case, and
Chapter~\ref{ch:state} adds a state diagram.

%============================================================================
\chapter{Functional Modelling: Use Cases}
\label{ch:use-cases}

\section{Actors}
\label{sec:actors}

An actor is a role played by a person or by another system that interacts
with the system being analysed. The three human roles are the ones Task~2
already identified in its Table~2-1; the two system roles are the sources
the system reads.

\begin{table}[H]
  \centering
  \caption{Actors and the role each plays.}
  \label{tab:actors}
  % Fixed widths: "Reviewer or auditor" and "Trajectory data store" are both
  % too wide for an "l" column, and an "l" column cannot wrap, so the table
  % would run past the 16cm text block.
  \setlength{\tabcolsep}{4.5pt}
  \begin{tabular}{@{}>{\raggedright\arraybackslash}p{3.2cm}
                    >{\raggedright\arraybackslash}p{2.5cm}
                    >{\raggedright\arraybackslash}p{9.0cm}@{}}
    \toprule[1.5pt]
    \textbf{Actor} & \textbf{Kind} & \textbf{Role} \\
    \midrule
    Analyst & Human, primary &
      A person with a question about recorded movement who cannot write a
      data query. Starts almost every use case. \src{Task~2, Table~2-1} \\
    \addlinespace
    Reviewer or auditor & Human, secondary &
      Checks how a past answer was produced and re-runs it. \\
    \addlinespace
    Assigned customer & Human, secondary &
      The team assigned by the Moodle group table; confirms that the
      analysis matches the requirement. \src{Task~2, Part~A source~1} \\
    \addlinespace
    Trajectory data store & System &
      Holds the recorded movement the system reads. Never written to. \\
    \addlinespace
    Recorded benchmark & System &
      Holds published question-and-answer pairs used to measure the
      system's own accuracy. \\
    \bottomrule[1.5pt]
  \end{tabular}
\end{table}

\tbd{whether the assigned customer team introduces an actor of its own ---
the actor list cannot be closed until Part~A source~1 has been elicited}

\section{Use Case Diagram}
\label{sec:uc-diagram}

Figure~\ref{fig:uc} shows the boundary of the system and the use cases it
offers. The six use cases correspond one to one with the six functional
requirements of Task~2, so the diagram can be read against Task~2 without a
translation table.

\begin{figure}[H]
  \centering
  \diagram{uc-diagram}{Use case diagram: the actors, the system boundary and
    the six use cases. Source: \texttt{figures/uc-diagram.puml}.}
  \caption{Use case diagram.}
  \label{fig:uc}
\end{figure}

Four of the six use cases are refinements of one another rather than
independent goals, and the diagram records this with \texttt{include}: a
question cannot be answered without first being interpreted, and a figure
cannot be drawn without first retrieving the records. \texttt{UC-5} and
\texttt{UC-6} stand on their own, because the accuracy dashboard and the
re-run of a past exchange can each be used without asking a new question.

\section{Use Case Descriptions}
\label{sec:uc-desc}

Each use case is described below by the same seven fields, so that the
descriptions can be compared with one another and, later, with the sequence
diagrams of Chapter~\ref{ch:seq}.

\subsection{UC-1: Ask a Question in Ordinary Language}
\label{uc:1}
\ucdesc{Analyst}
  {Obtain an exact written request from a question typed in ordinary
   language, without writing a query.}
  {The analyst has opened the system. The trajectories of interest are
   already held in the data store.}
  {\tbd{the numbered steps, once the format has been agreed with the
   customer. Draft from Task~2 REQ-1.1 and Task~1 \S3.4, feature F1.}}
  {\tbd{what happens when the question cannot be interpreted at all}}
  {The derived request is held for the rest of the session and shown back to
   the analyst beside the question.}
  {REQ-1.1}

\subsection{UC-2: Retrieve and Filter Trajectories}
\label{uc:2}
\ucdesc{System (on behalf of the Analyst)}
  {Find the trajectories the written request refers to and return only the
   matching subset.}
  {A written request exists.}
  {\tbd{the numbered steps. Draft from Task~2 REQ-1.2 and Task~1 \S3.4,
   feature F2.}}
  {\tbd{what happens when no record matches, and when a place name in the
   request cannot be resolved}}
  {A subset of the records is selected and identified by a snapshot
   identifier.}
  {REQ-1.2}

\subsection{UC-3: Select and Draw the Figure}
\label{uc:3}
\ucdesc{System (on behalf of the Analyst)}
  {Choose the chart that fits the question and render it, with no chart type
   or plotting tool chosen by the analyst.}
  {A subset of records has been selected.}
  {\tbd{the numbered steps. Draft from Task~2 REQ-1.3 and Task~1 \S3.4,
   feature F3.}}
  {\tbd{what happens when no chart type fits the question}}
  {A figure is produced and held with the exchange.}
  {REQ-1.3}

\subsection{UC-4: Summarise the Result in Words}
\label{uc:4}
\ucdesc{System (on behalf of the Analyst)}
  {Describe what the figure shows in one or two sentences, including the
   extremes found.}
  {A figure exists for this exchange.}
  {\tbd{the numbered steps. Draft from Task~2 REQ-1.4 and Task~1 \S3.4,
   feature F4.}}
  {\tbd{what the system says when the summary cannot be produced}}
  {A sentence or two is shown beside the figure.}
  {REQ-1.4}

\subsection{UC-5: Show the Accuracy Dashboard}
\label{uc:5}
\ucdesc{Reviewer or auditor}
  {See how accurate the automatic interpretation currently is, measured
   against the recorded benchmark.}
  {The benchmark is available.}
  {\tbd{the numbered steps. Draft from Task~2 REQ-1.5 and Task~1 \S3.4,
   feature F5. \textbf{The accuracy figure itself is still open in Task~2
   and must not be invented here.}}}
  {\tbd{what happens when the benchmark cannot be reached}}
  {A measured accuracy is displayed beside the answers.}
  {REQ-1.5}

\subsection{UC-6: Record and Re-run an Exchange}
\label{uc:6}
\ucdesc{Reviewer or auditor}
  {See how a past answer was produced, and reproduce the same figure from
   the record alone.}
  {At least one exchange has been made.}
  {\tbd{the numbered steps. Draft from Task~2 REQ-1.6 and Task~1 \S3.4,
   feature F6.}}
  {\tbd{what happens when the data behind a record has since changed}}
  {The same figure is rebuilt from the stored question, request, snapshot
   identifier and chart choice.}
  {REQ-1.6}

\section{Traceability to the Requirements}
\label{sec:uc-trace}

\begin{table}[H]
  \centering
  \caption{Every functional requirement of Task~2 is carried by exactly one
    use case.}
  \label{tab:uc-trace}
  % Fixed widths: an "l" description column would be widened by the longest
  % cell and would push the table past the text block.
  \setlength{\tabcolsep}{4.5pt}
  \begin{tabular}{@{}l l >{\raggedright\arraybackslash}p{9.0cm}@{}}
    \toprule[1.5pt]
    \textbf{Requirement} & \textbf{Use case} & \textbf{What the use case
      adds to the requirement} \\
    \midrule
    REQ-1.1 & UC-1 & Names the actor, the precondition and the failure
      case that the requirement states only in outline. \\
    \addlinespace
    REQ-1.2 & UC-2 & Fixes the snapshot identifier as the handle by which
      later steps refer to the selected subset. \\
    \addlinespace
    REQ-1.3 & UC-3 & Records that chart choice is the system's, not the
      analyst's. \\
    \addlinespace
    REQ-1.4 & UC-4 & Makes the summary a precondition of a finished
      exchange rather than an optional extra. \\
    \addlinespace
    REQ-1.5 & UC-5 & Separates the accuracy dashboard from the question
      round trip, which the requirement does not do. \\
    \addlinespace
    REQ-1.6 & UC-6 & States that re-running needs the record alone, not the
      data as it happens to stand now. \\
    \bottomrule[1.5pt]
  \end{tabular}
\end{table}

%============================================================================
\chapter{Static Modelling: Class and Object Diagrams}
\label{ch:static}

\section{Finding the Classes}
\label{sec:finding-classes}

The classes were taken from the nouns the Task~2 requirements and use case
descriptions use to talk about the domain, keeping those the system must
remember and dropping those that are only a way of speaking. Each class
below is given a name, the attributes it holds and the operations it
offers; a class with attributes but no operations would be a data record,
not a class, and the analysis is not complete until the operations have
been decided. \src{SE Chapter~5, object-oriented analysis}

\begin{table}[H]
  \centering
  \caption{Candidate classes and the reason each is kept or dropped.}
  \label{tab:classes}
  \setlength{\tabcolsep}{4.5pt}
  \begin{tabular}{@{}>{\raggedright\arraybackslash}p{2.8cm}
                    >{\raggedright\arraybackslash}p{5.2cm}
                    >{\raggedright\arraybackslash}p{6.7cm}@{}}
    \toprule[1.5pt]
    \textbf{Class} & \textbf{Where it comes from} & \textbf{Why it is a
      class} \\
    \midrule
    C-1 Trajectory & Movement records & The system must remember it, and it
      is the subject of almost every question. \\
    \addlinespace
    C-2 Waypoint & Positions along a trajectory & A position has no meaning
      apart from the trajectory that contains it, which makes the
      relationship a composition. \\
    \addlinespace
    C-3 MapLayer & Background map data & Layers are held and switched
      independently of any one trajectory. \\
    \addlinespace
    C-4 POI & Named places & Attached to a layer; a place can be
      interesting to several layers. \\
    \addlinespace
    C-5 Chart & Figures the system draws & Generalises the three chart
      kinds, which differ in how they are drawn but share a title, a scale
      and an axis convention. \\
    \addlinespace
    C-6 Query & The analyst's question & Must be remembered so that an
      exchange can be re-run. \\
    \addlinespace
    C-7 Exchange & One question and its answer & The unit the record and
      the re-run of UC-6 operate on. \\
    \addlinespace
    C-8 Benchmark & Published question-and-answer pairs & Read, never
      written; carried as a class so that UC-5 has something to run
      against. \\
    \addlinespace
    C-9 AccuracyReport & The measured accuracy & What UC-5 displays. \\
    \bottomrule[1.5pt]
  \end{tabular}
\end{table}

\tbd{the classes that the assigned customer's own data model implies, once
Part~A source~1 has been elicited}

\section{Class Diagram}
\label{sec:class-diagram}

Figures~\ref{fig:class} and~\ref{fig:class-ex} show the classes with their
attributes and their operations, and the relationships between them. They
are drawn as two diagrams rather than one: a single diagram holding all
twelve classes comes out about 1900 pixels wide, which at the width of the
text block scales the labels below 6~pt and makes them unreadable, whereas
each half stays above 9~pt. Figure~\ref{fig:class} holds the data classes
and Figure~\ref{fig:class-ex} the classes that carry a question through to
a recorded answer.

\begin{figure}[H]
  \centering
  \diagram{class-diagram-domain}{Class diagram, part 1: the data classes.
    Source: \texttt{figures/class-diagram-domain.puml}.}
  \caption{Class diagram, part 1: the data classes.}
  \label{fig:class}
\end{figure}

\begin{figure}[H]
  \centering
  \diagram{class-diagram-exchange}{Class diagram, part 2: the
    question-and-answer classes. Source:
    \texttt{figures/class-diagram-exchange.puml}.}
  \caption{Class diagram, part 2: the question-and-answer classes.}
  \label{fig:class-ex}
\end{figure}

\section{Relationships}
\label{sec:relationships}

\begin{table}[H]
  \centering
  \caption{Relationships between the classes, read from
    Figures~\ref{fig:class} and~\ref{fig:class-ex}.}
  \label{tab:relationships}
  % "TrajectoryMap, ODFlowChart, HeatMap" is far too wide for an "l" column
  % and cannot wrap in one, so every column here is given a fixed width.
  \setlength{\tabcolsep}{4.5pt}
  % The first and third columns are sized by their longest single word,
  % "AccuracyReport" and "Generalisation": a p{} column wraps between words
  % but a word wider than the column still runs into the margin.
  \begin{tabular}{@{}>{\raggedright\arraybackslash}p{2.85cm}
                    >{\raggedright\arraybackslash}p{3.9cm}
                    >{\raggedright\arraybackslash}p{2.7cm}
                    >{\raggedright\arraybackslash}p{5.25cm}@{}}
    \toprule[1.5pt]
    \textbf{Whole} & \textbf{Part} & \textbf{Kind} & \textbf{Why} \\
    \midrule
    Trajectory & Waypoint & Composition & A waypoint has no meaning apart
      from its trajectory; deleting the trajectory deletes its waypoints. \\
    \addlinespace
    MapLayer & POI & Aggregation & A place can outlive the layer it is
      currently shown on. \\
    \addlinespace
    Chart & TrajectoryMap, ODFlowChart, HeatMap & Generalisation & The
      three differ only in how they draw; the shared parts belong in the
      parent. \\
    \addlinespace
    Exchange & Query, Chart & Aggregation & An exchange is built from a
      question and the figure that answered it. \\
    \addlinespace
    AccuracyReport & Benchmark & Association & An accuracy figure is only
      meaningful against the benchmark it was measured on. \\
    \bottomrule[1.5pt]
  \end{tabular}
\end{table}

\section{Object Diagram}
\label{sec:object-diagram}

Part~A of the task description names the object diagram as the static model
to be provided. A class diagram says what kinds of thing may exist; an
object diagram says what exists at one moment, with the links between the
individual objects shown. Figure~\ref{fig:object} is the state of the
system immediately after one question has been answered, which is the
moment that makes every class in Figures~\ref{fig:class}
and~\ref{fig:class-ex} visible at once.

\begin{figure}[H]
  \centering
  \diagram{object-diagram}{Object diagram: one answered exchange, with the
    links between the individual objects. Source:
    \texttt{figures/object-diagram.puml}.}
  \caption{Object diagram for one answered exchange.}
  \label{fig:object}
\end{figure}

The two diagrams are consistent by construction: every object in
Figure~\ref{fig:object} is an instance of exactly one class in
Figures~\ref{fig:class} and~\ref{fig:class-ex}, every link between objects
is an instance of one
relationship in Table~\ref{tab:relationships}, and no object appears that
has no class.

\section{Traceability to the Requirements}
\label{sec:static-trace}

\begin{table}[H]
  \centering
  \caption{Classes, and the requirements they carry out.}
  \label{tab:class-trace}
  \setlength{\tabcolsep}{4.5pt}
  \begin{tabular}{@{}>{\raggedright\arraybackslash}p{3.8cm}
                    >{\raggedright\arraybackslash}p{3.0cm}
                    >{\raggedright\arraybackslash}p{7.9cm}@{}}
    \toprule[1.5pt]
    \textbf{Class} & \textbf{Requirement} & \textbf{Operation that carries
      it out} \\
    \midrule
    C-6 Query & REQ-1.1 & Interpreting a question written in ordinary
      language into an exact request. \\
    \addlinespace
    C-1 Trajectory & REQ-1.2 & Filtering the stored trajectories against a
      written request. \\
    \addlinespace
    C-5 Chart & REQ-1.3 & Choosing the chart kind that fits the question and
      rendering it. \\
    \addlinespace
    C-7 Exchange & REQ-1.4, REQ-1.6 & Holding the summary with the figure,
      and re-running the exchange from the record alone. \\
    \addlinespace
    C-8 Benchmark, C-9 AccuracyReport & REQ-1.5 & Re-running the recorded
      benchmark and reporting the measured accuracy. \\
    \bottomrule[1.5pt]
  \end{tabular}
\end{table}

%============================================================================
% Part A of the task description asks for a sequence diagram for *each* use
% case, so there is one section per use case here; that is what fixes the
% number of use cases in Chapter 2.
%============================================================================
\chapter{Dynamic Modelling: Sequence Diagrams}
\label{ch:seq}

This chapter gives one sequence diagram for each of the six use cases of
Chapter~\ref{ch:use-cases}, as Part~A of the task description requires. A
sequence diagram shows the messages that pass between the participants
while the use case runs, in the order in which they pass, which is the
dynamic counterpart of the class diagram: the classes say who may exchange
messages, the sequence diagrams say in what order they do.

\begin{figure}[H]
  \centering
  \diagram{seq-uc1}{Sequence diagram for UC-1, asking a question in
    ordinary language.}
  \caption{Sequence diagram for UC-1.}
  \label{fig:seq1}
\end{figure}

\begin{figure}[H]
  \centering
  \diagram{seq-uc2}{Sequence diagram for UC-2, retrieving and filtering
    trajectories.}
  \caption{Sequence diagram for UC-2.}
  \label{fig:seq2}
\end{figure}

\begin{figure}[H]
  \centering
  \diagram{seq-uc3}{Sequence diagram for UC-3, selecting and drawing the
    figure.}
  \caption{Sequence diagram for UC-3.}
  \label{fig:seq3}
\end{figure}

\begin{figure}[H]
  \centering
  \diagram{seq-uc4}{Sequence diagram for UC-4, summarising the result in
    words.}
  \caption{Sequence diagram for UC-4.}
  \label{fig:seq4}
\end{figure}

\begin{figure}[H]
  \centering
  \diagram{seq-uc5}{Sequence diagram for UC-5, showing the accuracy
    dashboard.}
  \caption{Sequence diagram for UC-5.}
  \label{fig:seq5}
\end{figure}

\begin{figure}[H]
  \centering
  \diagram{seq-uc6}{Sequence diagram for UC-6, recording and re-running an
    exchange.}
  \caption{Sequence diagram for UC-6.}
  \label{fig:seq6}
\end{figure}

\section{Consistency with the Class Diagram}
\label{sec:seq-consistency}

Every message in Figures~\ref{fig:seq1} to~\ref{fig:seq6} is sent to a
lifeline whose type is a class in Figures~\ref{fig:class}
and~\ref{fig:class-ex}, and every message name is an operation of that
class. Conversely, every operation in Figures~\ref{fig:class}
and~\ref{fig:class-ex} appears in at least one sequence diagram; an
operation that no use case calls would be a class member with no
requirement behind it.

%============================================================================
% Task description Part A(4): further diagrams are welcome.  The state
% diagram is the one extra that earns its place, because exactly one object
% in the model --- the exchange --- has behaviour that depends on what state
% it is in rather than only on the messages it receives.
%============================================================================
\chapter{Behavioural Modelling: State Diagram}
\label{ch:state}

Part~A(4) of the task description welcomes further diagrams beyond the three
it requires. The one added here is a state diagram, because the model
contains exactly one object whose behaviour depends on the state it is in
rather than only on the messages it receives: an exchange, which Task~2's
REQ-1.6 requires to be re-runnable after it is finished, and which may be
re-run more than once.

\begin{figure}[H]
  \centering
  \diagram{state-exchange}{State diagram for one exchange, from being asked
    to being re-run.}
  \caption{State diagram for an exchange.}
  \label{fig:state}
\end{figure}

The transition from the finished state back to the re-run state is what
makes the exchange re-runnable rather than a record that can only be read.
\tbd{any further state-dependent object the assigned customer's own
processes imply, once Part~A source~1 has been elicited}

%============================================================================
% \begin{thebibliography} is itself a \chapter*, so the \addcontentsline
% must sit inside it; placed outside, the contents entry would take the
% page number of the previous chapter's last page.
%============================================================================
\begin{thebibliography}{9}
\addcontentsline{toc}{chapter}{References}
\markboth{References}{}

\bibitem{task1} Group~1, \emph{TrajVis Studio: A Self-Service Platform for
  Natural-Language Trajectory Analysis}, Task~1 project proposal, Software
  Engineering, MUST, 2026.

\bibitem{task2} Group~1, \emph{TrajVis Studio: User Requirements
  Specification}, Task~2 deliverable, Software Engineering, MUST, 2026.

\bibitem{task3desc} \emph{Requirements Analysis}, Software Engineering,
  MUST, 2026.

\bibitem{template} \emph{Surface for Requirements Analysis}, template
  supplied with Task~3, Software Engineering, MUST, 2026.

\bibitem{repo} Project repository,
  \url{https://github.com/chenhongzhou666/MUST-SE-TrajVis}

\bibitem{ch5} \emph{SE Chapter 5: Requirements Analysis II (Object-Oriented
  Analysis)}, lecture slides, Software Engineering, MUST, 2026.09.30.

\end{thebibliography}

\end{document}
